\documentclass[../../../../SoftwareRequirementsSpecification.tex]{subfiles}
\begin{document}
% Richiede le macro definite in src/macros/useCaseMacros.tex
% (incluse nel preambolo di SoftwareRequirementsSpecification.tex)

\ucstart
\begin{xltabular}{\textwidth}{|l|X|}
  \hline
  \ucid{U-01} \\ \hline
  \ucname{Login} \\ \hline
  \ucactor{Utente (PT o Atleta)} \\ \hline
  \ucdescription{L'utente accede al sistema con la propria email e la
    propria password.} \\ \hline
  \ucprecond{L'utente deve essere registrato nel sistema.} \\ \hline
  \ucpostcond{L'utente è autenticato e può usare le funzioni previste
    per il proprio ruolo.} \\ \hline
  \ucmainflow{
    \item L'utente apre la pagina di login.
    \item Il sistema mostra il form di inserimento di email e
      password.
    \item L'utente inserisce la propria email e la propria password.
    \item L'utente preme il pulsante "Accedi".
    \item Il sistema verifica le credenziali confrontando l'email con
      il database e la password con il valore memorizzato (vedi le
      Note).
    \item Il sistema autentica l'utente e lo indirizza alla propria
      area.
  } \\ \hline
  \ucaltflow{
    \textbf{5a. Credenziali non valide:} \newline
    - L'email non è registrata, oppure la password non corrisponde. \newline
    - Il sistema mostra un unico messaggio di errore ("Email o
    password non validi"), senza distinguere i due casi (vedi le
    Note). \newline
    - Il flusso riprende dal passo 2. \newline
    \textbf{6a. Primo accesso dell'Atleta:} \newline
    - L'utente autenticato è un Atleta che sta accedendo con la
    password temporanea generata dal sistema in PT-02, e non ha
    ancora scelto una propria password né inserito i propri dati
    anagrafici e fisici. \newline
    - Il sistema lo indirizza al caso d'uso A-01 (Completamento
    Profilo) prima di dargli accesso alle altre funzioni.
  } \\ \hline
  \ucnote{
    Il sistema non memorizza la password in chiaro, ma solo la sua
    impronta calcolata con una funzione di hash. La verifica del
    passo 5 confronta le impronte, non le password. \newline
    Il messaggio di errore del flusso 5a è volutamente unico: se il
    sistema distinguesse "email non registrata" da "password errata",
    permetterebbe a chiunque di scoprire quali indirizzi email sono
    registrati nel sistema. \newline
    L'autenticazione è uguale per entrambi gli attori. Ciò che
    cambia, dopo il login, sono le funzioni a cui l'utente ha
    accesso.
  } \\ \hline
\end{xltabular}
\end{document}
